\documentclass[11pt,a4paper]{article}
\usepackage[T1]{fontenc}
\usepackage{lmodern,amsmath,amssymb,amsthm,microtype,booktabs}
\usepackage[margin=25mm]{geometry}
\usepackage[colorlinks=true,urlcolor=blue,linkcolor=blue]{hyperref}
\newcommand{\dd}{\mathop{}\!\mathrm d}
\newcommand{\Id}{\mathbf 1}
\title{Sewing measures and identity limits\\of the recoil two-point spectrum}
\author{Alan Day \and Chengheng\thanks{AI research assistant (OpenAI).}}
\date{Working research note --- 31 October 2026}
\begin{document}\maketitle
\begin{abstract}
We distinguish the identity representation in real-parameter Virasoro
TQFT from a normalized massless endpoint of the recoil two-point
spectrum. The common cone-angle continuation used in the paired
amplitude cannot reach the vacuum of both Virasoro factors.
The reduced spectrum nevertheless admits a controlled delta-function
limit. We determine its pole-crossing correction and a finite-mass
error bound. The previously exhibited insertion freedom survives
this endpoint test. These statements do not construct a de Sitter
sewing inner product or a physical finite-gravity four-point function.
\end{abstract}

\section{What the known sewing rule fixes}
We use CEZ \cite{cez}, version 3, whose HTML equations have continuous
numbering. In the standard block normalization, for real \(b>0\)
and positive real internal momenta, equations (23)--(25) give
\begin{align}
 \langle s|s'\rangle&=\frac{\delta(s-s')}{\rho_s(s)},\\
 \rho_s(s)&=\rho_0(s)C_0(P_1,P_2,s)C_0(P_3,P_4,s),\\
 \rho_0(P)&=4\sqrt2\sinh(2\pi bP)\sinh(2\pi P/b).
\end{align}
Thus the four-punctured-sphere resolution is
\(\Id=\int_0^\infty\dd s\,\rho_s(s)|s\rangle\langle s|\).
Multiplying each block by the square roots of its vertex
\(C_0\) factors changes the measure to \(\rho_0\), as in CEZ
(27)--(28). Dropping the vertices without this basis change is
inconsistent.

With the convention \(|s\rangle=\int_0^\infty\dd t\,F_{st}|t\rangle\),
preservation of the inner product requires
\begin{equation}
 \int_0^\infty\dd t\,\frac{F_{st}\overline{F_{s't}}}{\rho_t(t)}
 =\frac{\delta(s-s')}{\rho_s(s)}. \label{eq:unitarity}
\end{equation}
This follows from the displayed block metric; CEZ (55) and (122)
supply the corresponding exchange and fusion-inverse relations.
The bilinear inverse relation and the sesquilinear isometry
must not be confused. We have not evaluated the full fusion kernel
numerically in this note.

A change \(|s\rangle\mapsto a_s|s\rangle\), with nonzero \(a_s\),
leaves \eqref{eq:unitarity} covariant if
\begin{equation}
 \rho'_s=\rho_s/|a_s|^2,\qquad F'_{st}=a_sF_{st}/a_t.
\end{equation}
This is basis freedom. It does not show that an arbitrary change
of a physical operator insertion preserves crossing with a fixed
fusion kernel.

The identity limit in CEZ (29) is distributional:
\begin{equation}
 \lim_{P_3\to iQ/2}C_0(P_1,P_2,P_3)
 =\frac{\delta(P_1-P_2)}{\rho_0(P_1)},\qquad Q=b+b^{-1}.
 \label{eq:cezidentity}
\end{equation}
The momenta \(P_1,P_2>0\) remove reflection redundancy; the
full-line convention has an additional reflected delta.
The vacuum is outside the normalizable continuum. CEZ (46)
distinguishes an incoming or external identity from an outgoing
continuum label. In particular, squaring
\eqref{eq:cezidentity} pointwise is not a definition of a paired
amplitude: \(\delta^2\) is undefined.

\section{Why this is not yet the de Sitter measure}
CLP \cite{clp}, section 6, distinguishes analytic continuation of
fixed-topology amplitudes from a de Sitter theory with arbitrary
sewing. CEM \cite{cem}, equations (20), (24) and section 5,
uses a late-time Teichm\"uller/moduli-space inner product which
ordinary conformal blocks do not diagonalize.
These are statements about the versions inspected here, not an
exhaustive claim about all subsequent literature.

For a nonorthogonal basis, the required formal resolution is instead
\begin{equation}
 \Id=\iint\dd a\,\dd b\ |a\rangle(G^{-1})_{ab}\langle b|,
 \qquad G_{ab}=\langle a|b\rangle. \label{eq:gram}
\end{equation}
Existence of the inverse on a suitable test space, treatment of
null directions, positivity and convergence are additional tasks.
Neither the analytically continued Cardy factor nor the observer
weight \(z(m)=m^2e^{-2\pi m}\) has thereby been identified with the
physical Plancherel measure.

\section{The common-angle path cannot reach two vacua}
Use exactly the two factors in CLP (3.2), (3.4), (3.6) and (3.44):
\begin{equation}
 b_\pm^2=\pm i/\kappa,\qquad
 P_\pm=\pm\frac{i\alpha}{4\pi b_\pm},\qquad
 h_\pm=\frac{(b_\pm+b_\pm^{-1})^2}{4}+P_\pm^2.
\end{equation}
For a common angle \(\alpha=2\pi(1-\nu/\kappa)\), direct algebra gives
\begin{align}
 h_+&=\frac{1-i\nu}{2}+\frac{i(1+\nu^2)}{4\kappa},\\
 h_-&=\frac{1+i\nu}{2}-\frac{i(1+\nu^2)}{4\kappa},&
 h_++h_-&=1. \label{eq:weights}
\end{align}
This exact label identity is consistent with CLP's framing discussion.
At \(\nu=i\), the pair is \((1,0)\); at \(\nu=-i\), it is \((0,1)\).
At \(\nu=0\) neither factor is the vacuum.
Here \(h_\pm\) label the two continued factors; their sum is not
being identified with the physical scalar dimension.

On the positive-angle branch, the separate vacuum conditions require
\begin{equation}
 \alpha^{\rm vac}_\pm=2\pi(1+b_\pm^2)
 =2\pi(1\pm i/\kappa). \label{eq:vacangles}
\end{equation}
They are distinct at finite \(\kappa\).
The reflected negative-angle solutions do not restore a common angle.
Although these angles approach each other as \(\kappa\to\infty\),
their difference is of the order retained by the fixed-mass
scaling. It cannot be discarded when testing a vacuum limit.

Consequently, the scalar relation \(m_{\rm field}^2=1+\nu^2\)
does not justify identifying \(\nu=\pm i\) with a double-vacuum
leg of the paired TQFT graph. Continuing the real-\(\nu\) paired
expression meromorphically with a common complex \(\nu\) is also
different from taking a fresh modulus square at each complex \(\nu\).
A double-vacuum calculation needs independently specified labels
and a pairing prescription.

\section{A controlled endpoint of the reduced spectrum}
We now change the question explicitly. Work at radius one,
after the fixed-mass \(\kappa\to\infty\) limit, and continue only
the established reduced two-sided spectrum
\begin{equation}
 k_\nu(\omega)=\frac1{8\pi[\cosh\pi\omega+\cosh\pi\nu]},
 \qquad g_\nu(0)=\frac{\nu}{4\pi\sinh\pi\nu}.
\end{equation}
Set \(\nu=i(1-\epsilon)\), \(0<\epsilon<1\).
Dividing by the divergent
\(g_\nu(0)=(1-\epsilon)/(4\pi\sin\pi\epsilon)\) gives
\begin{align}
 p_\epsilon(\omega)&=
 \frac{\sin\pi\epsilon}
 {2(1-\epsilon)[\cosh\pi\omega-\cos\pi\epsilon]},\\
 \chi_\epsilon(t)&=\int_{\mathbb R}\dd\omega\,
 e^{-i\omega t}p_\epsilon(\omega)
 =\frac{\sinh((1-\epsilon)t)}{(1-\epsilon)\sinh t}.
 \label{eq:char}
\end{align}
At \(t=0\) the continuous value is one. No \(1/(2\pi)\) occurs
in this forward transform. The identity follows by analytic
continuation inside \(0<\operatorname{Im}\nu<1\), where compact
subsets admit integrable spectral domination, or by residues.

The density is nonnegative and integrates to one. Differentiating
\eqref{eq:char} at zero gives
\begin{equation}
 v_\epsilon:=\int_{\mathbb R}\omega^2p_\epsilon(\omega)\dd\omega
 =\frac{2\epsilon-\epsilon^2}{3},\qquad
 \int_{|\omega|>D}p_\epsilon(\omega)\dd\omega\le\frac{v_\epsilon}{D^2}.
 \label{eq:moment}
\end{equation}
Hence \(p_\epsilon\rightharpoonup\delta\).
For a \(C^2\) test function with bounded second derivative,
Taylor's formula and evenness give the explicit error
\(\|f''\|_\infty v_\epsilon/2\).
Concentration and normalization also prove weak convergence for
bounded continuous tests. This is a normalized infrared endpoint,
not a finite unnormalized massless-field correlator or a
construction of the identity operator on a gravitational Hilbert space.

\subsection*{The pole crossing is visible}
The poles are \(\omega=i(2n\pm\epsilon)\).
The nearest two, \(\pm i\epsilon\), pinch the real axis at zero.
Their residues are respectively
\(\pm[2\pi i(1-\epsilon)]^{-1}\).
No pole crosses before the endpoint on the positive-\(\epsilon\)
path. Continuing through it is a different matter.

For \(\epsilon=-\delta\), \(0<\delta<1\), the same real-axis integral is
\begin{equation}
 \int_{\mathbb R}\dd\omega\,e^{-i\omega t}p_{-\delta}(\omega)
 =-\frac{\sinh((1-\delta)t)}{(1+\delta)\sinh t}.
\end{equation}
It differs from the continued right side of \eqref{eq:char} by
\begin{equation}
 \chi_{-\delta}(t)
 -\int_{\mathbb R}\dd\omega\,e^{-i\omega t}p_{-\delta}(\omega)
 =\frac{2\cosh(\delta t)}{1+\delta}. \label{eq:residues}
\end{equation}
This is precisely the sum of the two crossing-pole contributions;
it also follows from the hyperbolic addition formula.
The difference survives at \(t=0\).
We have tracked this reduced two-pole problem, not all poles of
the finite-\(\kappa\) double-sine kernel.

\section{Finite observer mass does not remove the freedom}
Fix \(M>m_*>0\), \(D=2(M-m_*)\), \(A=\cosh(2\pi M)\). Define
\begin{align}
 R_{M,\epsilon}(\omega)&=
 \left(1-\frac{\omega^2}{4M^2}\right)
 \frac{A-\cosh\pi\omega}{A-\cos\pi\epsilon},\\
 q_{M,\epsilon,\lambda}(\omega)&=
 \Id_{[-D,D]}(\omega)p_\epsilon(\omega)R_{M,\epsilon}(\omega)
 \left(1+\frac{\lambda\omega^2}{4M^2}\right),\quad\lambda\ge0 .
 \label{eq:q}
\end{align}
This is a declared spectral model, with \(0\le R\le1\) on its support.
Let \(T_D=\int_{|\omega|>D}p_\epsilon\) and
\(v_D=\int_{-D}^D\omega^2p_\epsilon\). The exact mass loss is
\begin{equation}
 1-\int q_{M,\epsilon,0}
 =T_D+\frac{v_D}{4M^2}
 +\frac{\sin\pi\epsilon}{(1-\epsilon)(A-\cos\pi\epsilon)}
 \left(D-\frac{D^3}{12M^2}\right). \label{eq:loss}
\end{equation}
To prove it, expand \(1-R\). In its second term,
\(\cosh\pi\omega-\cos\pi\epsilon\) cancels the spectral denominator,
leaving a quadratic polynomial to integrate.
Using \eqref{eq:moment}, set
\begin{equation}
 B_{M,\epsilon}=
 \frac{v_\epsilon}{D^2}+\frac{v_\epsilon}{4M^2}
 +\frac{\sin\pi\epsilon}{(1-\epsilon)(A-\cos\pi\epsilon)}
 \left(D-\frac{D^3}{12M^2}\right)=O_{M,m_*}(\epsilon).
\end{equation}
For all real \(t\), the triangle inequality gives
\begin{equation}
 |\widehat q_{M,\epsilon,\lambda}(t)-\chi_\epsilon(t)|
 \le B_{M,\epsilon}+\frac{\lambda v_\epsilon}{4M^2}.
 \label{eq:bound}
\end{equation}
The same bound controls \(\|q_{M,\epsilon,\lambda}-p_\epsilon\|_1\).
Therefore every fixed \(\lambda\ge0\) yields the same weak limit
\(q\rightharpoonup\delta\); restoring the common factor gives
\(z(M)\delta\). Yet these same choices change the first finite-mass
correction on the real principal series. Leading QFT matching plus
this normalized endpoint test thus fails to fix the insertion.

More generally, a fixed bounded continuous even factor, or a family
converging uniformly to one on this finite interval, is tested only
through its limiting diagonal value \(n(M,0)\). This is a necessary
normalization condition in the chosen spectral convention, not
a proof that all such factors extend to a crossing-consistent
physical four-point function.

\section{Limits, verification and the next physical task}
The identity limit is not uniform over time:
\begin{equation}
 \lim_{\epsilon\downarrow0}\chi_\epsilon(t)=1
 \quad(t\ {\rm fixed}),\qquad
 \lim_{|t|\to\infty}\chi_\epsilon(t)=0
 \quad(\epsilon>0),
\end{equation}
and for fixed \(\tau\ne0\),
\begin{equation}
 \lim_{\epsilon\downarrow0}\chi_\epsilon(\tau/\epsilon)
 =e^{-|\tau|}.
\end{equation}
Bound \eqref{eq:bound} transfers the latter result to \(q\) at fixed
\(M,m_*,\lambda\). This is an infrared model timescale, not a
finite-gravity chaos exponent.

The saved diagnostics contain 106 passed checks: 16 real-axis
Fourier/moment checks, 18 crossing checks, 36 finite-mass identity,
bound and precision comparisons, and 36 label-algebra checks.
The last group includes repeated vacuum checks across different
probe-label loops; it is not 36 independent physical tests.
Calculations used 50 decimal digits and selected finite-mass
integrals were repeated at 65 digits; the largest change was
\(6.41\times10^{-52}\). This is an empirical refinement, not
interval certification. Parameters, versions, outputs and hashes
are stored with the code. Large-time scaling entries evaluate
the exact transform; they are not claimed as oscillatory quadratures.

The next concrete task is the global \(\mathrm{SL}(2,\mathbb C)\)
partial-wave inner product in the \(G=0\) limit, including reflection
and spectral measure, followed by an explicit comparison with the
regularized recoil graph \(H\). A Barnes-integral resemblance is
insufficient. Any double-vacuum study must instead keep independent
labels and track the continuation of already-sewn expressions.
Physical finite-\(\kappa\) sewing, finite-\(G\) four-point dynamics
and a complete de Sitter dual remain open.

\begin{thebibliography}{9}
\bibitem{cez} S. Collier, L. Eberhardt and M. Zhang,
\emph{Solving 3d Gravity with Virasoro TQFT},
\href{https://arxiv.org/html/2304.13650v3}{arXiv:2304.13650v3}.
\bibitem{clp} J. Chandra, D. Liska and B. Post,
\emph{3D de Sitter from complex Virasoro TQFT: Observers and knots},
\href{https://arxiv.org/html/2610.30247v1}{arXiv:2610.30247v1}.
\bibitem{cem} S. Collier, L. Eberhardt and B. M\"uhlmann,
\emph{A microscopic realization of dS$_3$},
\href{https://arxiv.org/html/2501.01486v1}{arXiv:2501.01486v1}.
\end{thebibliography}
\end{document}
